% AgentSynth technical report.
% Build: tectonic main.tex   (CI does this on every push touching paper/)
% Every number here is regenerated by paper/experiments.py — run it before editing a table.
% Set the real author list before the arXiv submission.
\documentclass[11pt]{article}

\usepackage[margin=1.1in]{geometry}
\usepackage{amsmath}
\usepackage{booktabs}
\usepackage{microtype}
\usepackage[hidelinks]{hyperref}
\usepackage{xurl}
\usepackage{listings}
\usepackage{xcolor}

\lstset{
  basicstyle=\ttfamily\small,
  breaklines=true,
  frame=single,
  framerule=0.2pt,
  xleftmargin=6pt,
  aboveskip=8pt,
  belowskip=8pt,
}

\newcommand{\agentsynth}{\textsc{AgentSynth}}

\title{\agentsynth: Outcome-Verified Benchmarks and\\Reproducible Evaluation for LLM Agents}
\author{AgentSynth contributors\\[2pt]
  \small\url{https://github.com/agentsynth/agentsynth} \quad \url{https://agentsynth.tech}}
\date{July 2026 — v0.9 draft}

\begin{document}
\maketitle

\begin{abstract}
Agent benchmarks are graded on the wrong evidence. When the score comes from the
transcript --- did the model say it fixed the ticket, did an LLM judge like the
answer --- an agent can talk its way to a result, and recent audits show that the
major public benchmarks can be driven to near-perfect scores without solving a
single task. \agentsynth{} is an open toolchain built around a stricter rule:
\emph{a run passes only if the world it acted on ends in the goal state}.
Scenarios pair a live environment (a writable SQL database, a code sandbox, a
multi-turn user) with checkers that assert on that end state, so the transcript
is free to be wrong. Around this core the toolchain supplies the trust
infrastructure benchmarks usually lack: an adversarial audit that measures how
gameable a pack's checkers are before anyone trusts it; contamination tooling
(canaries, corpus overlap, held-out isomorphic siblings) that separates
``knows the answer'' from ``can do the work''; reliability reporting beyond
pass@1 (the full pass\textsuperscript{1}..pass\textsuperscript{k} decay with
confidence intervals); and reproducible run manifests whose hash any third party
can re-derive. Agents are evaluated as-is over stdio, HTTP, or MCP --- no
framework rewrite --- and every scenario doubles as an RL environment with a
verifiable reward, exportable to the \texttt{verifiers} and OpenEnv ecosystems.
We report the audit of our own shipped packs (robustness 86--100\%, including
the weaknesses the audit found in them), a sibling experiment in which an
answer-key replay collapses from 100\% to 0\% while an adaptive agent holds
100\%, and a reliability study in which a 77.5\% pass@1 agent is dependable on
only 30\% of tasks. Everything runs offline and every number in this report is
regenerated by one script in the repository.
\end{abstract}

\section{Introduction}

Two findings frame the current state of agent evaluation. First, scores on the
public agent benchmarks are inflated: contamination alone is estimated to add
5--15 points, and most leaderboards still report a single run of a stochastic
system. Second --- and more damning --- the benchmarks themselves are gameable.
A 2026 audit by Berkeley RDI constructed working exploits against eight
prominent agent benchmarks, achieving near-perfect scores with no task-solving
at all: a pytest hook that forces assertions to pass, a fake \texttt{curl} that
satisfies a network check, reading the benchmark's own answer configuration off
the local filesystem~\cite{rdi2026broke}. None of these attacks are clever
agent behavior; they are failures of \emph{verification design}. The scoring
functions trusted evidence the agent could fabricate.

The lesson is older than agents: an assertion is only as good as the state it
inspects. Test suites learned this long ago --- a test that greps a log is a
worse test than one that reads the database. Applied to agents, the rule is
that the benchmark, not the agent, must own the world, and the score must be a
function of that world's end state.

\agentsynth{} is a small open-source toolchain (Python, MIT,
\texttt{pip install agentsynth-ai}) that applies this rule end to end and then
takes the obvious next steps that follow from taking it seriously:

\begin{enumerate}
  \item \textbf{Outcome-verified scenarios} (\S\ref{sec:scenarios}). A scenario
  is a task over a live environment plus \emph{checkers} that assert on the
  environment's final state. The agent's own claims carry no score weight; even
  code is graded by hidden tests executed after the fact, not by inspection.
  \item \textbf{An adversarial audit for benchmarks}
  (\S\ref{sec:audit}). If checkers can be satisfied without doing the work, the
  benchmark is broken \emph{by construction}, so the toolchain ships
  knowledge-free adversaries (do nothing; say a canned phrase; echo the prompt;
  echo plus one probing call) and reports how much of a pack they defeat ---
  including packs we ship ourselves.
  \item \textbf{Contamination controls} (\S\ref{sec:contamination}). Canary
  strings, shingle-overlap scans against a training corpus, and \emph{held-out
  isomorphic siblings}: relabelled worlds where every relational fact still
  holds, so a memorized answer fails while a competent agent does not notice.
  \item \textbf{Reliability beyond pass@1} (\S\ref{sec:reliability}).
  $k$-trial benching reports the whole
  pass\textsuperscript{1}..pass\textsuperscript{k} decay curve with Wilson
  intervals and names the flaky scenarios, in the spirit of
  pass\^{}k~\cite{yao2024tau} and the argument of~\cite{beyondpass1}.
  \item \textbf{Reproducible run manifests} (\S\ref{sec:provenance}). Every
  bench emits a content hash of the pack, the policy, seeds, trials, and
  per-scenario outcomes folded into one \texttt{run\_hash}; a verifier re-runs
  the manifest and confirms the hash, and a tampered pack is detected by its
  fingerprint.
  \item \textbf{Evaluation without a rewrite} (\S\ref{sec:interfaces}). An
  agent is benched as-is: one JSON object per step over stdin/stdout or HTTP,
  or the pack served as an MCP server that any MCP client can work against.
  Scenarios export as \texttt{verifiers} environments and OpenEnv modules, so
  the same checkers that grade a benchmark also pay RL reward.
\end{enumerate}

The point of the toolchain is not another leaderboard. It is that
\emph{verification, audit, contamination control, and reproducibility are
properties a benchmark should carry with it}, the way a library carries its
tests. Nothing here requires our scenarios: the audit runs against any pack in
the schema, and the schema is deliberately small.

\section{Related work}\label{sec:related}

\textbf{Agent benchmarks.} $\tau$-bench and $\tau^2$-bench ground scoring in a
database state and introduced pass\^{}k reliability~\cite{yao2024tau,tau2bench};
SWE-bench grades patches by held-out tests~\cite{swebench}; WebArena, GAIA,
OSWorld, and Terminal-Bench broadened the task
surface~\cite{webarena,gaia,osworld}. \agentsynth{} generalizes the
state-graded pattern into a portable schema and adds the audit, contamination,
and provenance layers those benchmarks mostly lack --- the RDI exploits landed
on several of them precisely because the harness trusted agent-fabricable
evidence~\cite{rdi2026broke}.

\textbf{Evaluation harnesses.} HAL standardizes large-scale, cost-aware agent
evaluation with public logs~\cite{hal}; pytest-style checkers
(e.g.\ DeepEval) make LLM output assertions ergonomic; LiveBench refreshes
questions to resist contamination~\cite{livebench}. These grade model
\emph{output}; our checkers grade the \emph{world}, which is what makes the
knowledge-free audit meaningful.

\textbf{Benchmark integrity.} Berkeley RDI's exploit study and its BenchJack
scanner audit a benchmark's scoring pipeline from the attacker's
side~\cite{rdi2026broke}; adversarial hacker--fixer loops harden benchmarks
iteratively~\cite{hackerfixer}; reward-hacking of verifiers is now documented
at model scale~\cite{gamingverifiers}. \agentsynth{} is complementary
infrastructure: the audit is built into the pack lifecycle (a pack ships with
its robustness score), and the synthesis path (\S\ref{sec:synthesis}) derives
checkers that pass the audit by construction.

\textbf{RL environments.} The \texttt{verifiers} library and the Environments
Hub made RL environments a community artifact~\cite{verifiers}; OpenEnv is
emerging as the interoperability standard~\cite{openenv}. We do not compete
with either: packs export into both, carrying their audit and provenance
metadata with them.

\section{Outcome-verified scenarios}\label{sec:scenarios}

A \emph{scenario} is a declarative YAML/JSON record: a task string, an
environment specification, a step budget, and a list of checkers. The
environment is real --- a writable SQLite database built from the scenario's
schema and rows, a Python sandbox, a Docker container, a REST surface, a
headless browser, or any environment registered through an entry-point plugin.
The agent interacts through the environment's tools; every call executes.

\textbf{Checkers.} A checker is a small assertion over the environment's final
state: \texttt{SqlCheck} (a query whose result must equal / contain a value),
\texttt{CalledTool} (the episode used a required tool),
\texttt{AnswerContains} (the final answer names a fact --- used for read-only
and refusal tasks, and flagged as text-graded by the audit), and
\texttt{CodeCheck}, which collects the code the agent wrote during the episode
and executes a \emph{hidden} test against it in the sandbox. A scenario passes
only if every checker passes; the outcome score is the fraction that do.
Checkers run \emph{after} the episode on the same live environment, so a claim
in the transcript ("I refunded order 7") is worth nothing unless row 7 changed.

\textbf{Episodes.} \texttt{AgentGym} runs one episode: actions are tool calls
or a final answer, the step budget is enforced, and the terminal verdict
combines the outcome score with optional trajectory verification and an LLM
judge whose weight is zero in benchmark mode. Multi-turn tasks
(\texttt{run\_conversation}) replay scripted user turns against one persistent
environment and run checkers once at the end, in the style of
$\tau^2$-bench~\cite{tau2bench}: an agent that satisfies turn one but forgets
turn two fails.

\textbf{Packs and gates.} Scenarios ship in \emph{packs}. A pack must pass
\texttt{pack validate} before it ships: a reference oracle must solve every
scenario, and a do-nothing policy must solve none --- the two cheapest ways a
benchmark can be broken (unsolvable tasks; vacuous checkers) are gated out
mechanically. The repository ships four packs used throughout this report:
\texttt{core\_v1} (10 single-table business tasks), \texttt{core\_v2} (14
tiered tasks including multi-table consistency: refund \emph{and} restock,
void \emph{and} cancel), \texttt{policy\_v1} (tool use under an explicit
policy), and \texttt{code\_v1} (four coding tasks graded by hidden tests).

Table~\ref{tab:bench} shows the pass rates of four scripted policies. The
spread is the point: the oracle proves solvability at 100\%, the do-nothing
policy proves the checkers are not vacuous at 0\%, and the generic policies
land where a benchmark wants its headroom. Note \texttt{policy\_v1} and
\texttt{code\_v1} hold even a competent generic agent at 0\%: compliance with
a stated policy and passing hidden tests are not reachable by generic
tool-calling competence.

\begin{table}[t]
\centering
\begin{tabular}{lrrrr}
\toprule
pack & oracle & generic expert & read-only & do-nothing \\
\midrule
\texttt{core\_v1} (n=10) & 100\% & 10\% & 0\% & 0\% \\
\texttt{core\_v2} (n=14) & 100\% & 100\% & 29\% & 0\% \\
\texttt{policy\_v1} (n=4) & 100\% & 0\% & 0\% & 0\% \\
\texttt{code\_v1} (n=4) & 100\% & 0\% & 0\% & 0\% \\
\bottomrule
\end{tabular}
\caption{Scripted policies against the shipped packs (seed 7, offline). The
oracle--zero spread is a shipping gate: every scenario solvable, no checker
satisfiable by silence. The generic expert is a careful inspect-act-verify
loop with no pack-specific knowledge; \texttt{core\_v2} was tuned so that
careful generic competence clears it while carelessness does not.}
\label{tab:bench}
\end{table}

\section{The trust toolchain}

\subsection{Auditing a benchmark before trusting it}\label{sec:audit}

The RDI exploits worked because nobody attacked the benchmarks before the
models did. \texttt{agentsynth pack audit} runs that attack as a shipping
gate. Four \emph{knowledge-free} adversaries --- policies that know nothing
about the tasks --- are run over every scenario:

\begin{itemize}
  \item \textbf{noop}: end the episode immediately with an empty answer;
  \item \textbf{constant}: reply with one canned all-purpose sentence;
  \item \textbf{echo}: answer with the task text itself;
  \item \textbf{echo+probe}: make one generic tool call, then echo the task.
\end{itemize}

A scenario any of them passes is gameable by construction. The audit also
flags \emph{answer leaks} --- expected answer tokens that already sit in the
task prompt (matched on word boundaries) --- and scenarios whose state checks
are satisfiable by inaction. Table~\ref{tab:audit} is the audit of our own
packs, weaknesses included.

\begin{table}[t]
\centering
\begin{tabular}{lrrl}
\toprule
pack & robustness & gamed / n & what the audit found \\
\midrule
\texttt{core\_v1} & 90\% & 1/10 & refusal task falls to echo+probe; two leaks \\
\texttt{core\_v2} & 86\% & 2/14 & both refusal tasks fall to echo+probe \\
\texttt{policy\_v1} & 100\% & 0/4 & --- \\
\texttt{code\_v1} & 100\% & 0/4 & --- \\
\bottomrule
\end{tabular}
\caption{\texttt{pack audit} against the shipped packs. The failures are
informative: refusal tasks ("do not cancel order 9; explain why") are graded
partly on the answer text, and an adversary that echoes a prompt containing
the word ``shipped'' can pass that check. The audit labels these text-graded
scenarios so a consumer can weigh them; hardening options are to assert the
state \emph{change} instead, or to drop prompt words from the expected set.}
\label{tab:audit}
\end{table}

Two design choices matter. The adversaries are knowledge-free so the audit
measures the \emph{checkers}, not the difficulty of the tasks; and the audit
is honest about its own packs --- \texttt{core\_v2} ships with a documented
86\% floor rather than a hidden one. The audit also closes a loop with
synthesis: \texttt{pack new --from-demo} derives checkers by diffing the world
a demonstration produced (asserting the written value, keyed on primary keys,
with a row-count guard), which produces packs that pass the audit at 100\% by
construction.

\subsection{Contamination: knowing vs.\ doing}\label{sec:contamination}

Three mechanisms, in increasing order of strength:

\textbf{Canaries.} Every scenario mints a deterministic canary string
(\texttt{agentsynth-canary-\dots}); its verbatim appearance in a model output
or a training corpus is evidence the pack leaked.

\textbf{Corpus overlap.} Character-shingle overlap between each task and a
supplied corpus; in the calibration in Table~\ref{tab:siblings}, a task
planted verbatim in the corpus scores overlap 1.0 and is flagged while every
other scenario scores 0.0.

\textbf{Held-out isomorphic siblings.} \texttt{held\_out\_pack} relabels the
string constants of each scenario's world (names, e-mails, status labels) with
a seeded map, rewriting rows, task text, and checker constants consistently,
and preserving every number and relation. Crucially, values absent from the
initial world --- the values the agent is \emph{instructed to write}, like
\texttt{refunded} --- are untouched, so the sibling remains exactly as
solvable: same schema, same relational facts, different surface tokens.
What changes is only what memorization can exploit.

\begin{table}[t]
\centering
\begin{tabular}{lrr}
\toprule
agent & original & held-out sibling \\
\midrule
answer-key replay (id-keyed oracle) & 100\% & \textbf{0\%} \\
adaptive agent (reads the world; n=2 subset) & 100\% & \textbf{100\%} \\
\bottomrule
\end{tabular}
\caption{The sibling experiment on \texttt{core\_v1}. The replay agent
re-issues the oracle's memorized SQL for each scenario id --- the behavioral
signature of a contaminated model --- and collapses on siblings because the
memorized surface tokens no longer exist in the world. An adaptive agent that
reads the world and acts from the task's meaning does not notice the
relabelling. The gap between those two rows is the quantity a held-out
sibling measures.}
\label{tab:siblings}
\end{table}

\subsection{Reliability: the decay from pass@1 to pass\^{}k}\label{sec:reliability}

A single-run pass rate flatters a stochastic agent. \texttt{bench --trials k}
runs the pack $k$ times on shifted seeds and scores
pass\^{}k~\cite{yao2024tau}: a scenario counts only if \emph{every} trial
passed. The report includes the unbiased whole curve
pass\textsuperscript{1}..pass\textsuperscript{k} (via
$\binom{c}{k}/\binom{n}{k}$), a Wilson interval on each point, and the list of
flaky scenarios. Table~\ref{tab:reliability} shows the instrument on a
deliberately flaky solver --- the oracle, except a seeded coin makes it quit
early 25\% of the time, the shape of a real model's failure profile:

\begin{table}[t]
\centering
\begin{tabular}{lrrrr}
\toprule
 & pass\textsuperscript{1} & pass\textsuperscript{2} & pass\textsuperscript{3} & pass\textsuperscript{4} \\
\midrule
flaky solver on \texttt{core\_v1} & 77.5\% & 58.3\% & 42.5\% & \textbf{30.0\%} \\
\bottomrule
\end{tabular}
\caption{Reliability decay over 4 trials (95\% Wilson CI on
pass\textsuperscript{1}: 62.5--87.7\%; on pass\textsuperscript{4}:
10.8--60.3\%). An agent that looks like ``77.5\%'' in single-run reporting is
dependable on 30\% of tasks; 7 of 10 scenarios are flaky (pass sometimes).
Leaderboards built on this toolchain display the decay, not the flattering
endpoint.}
\label{tab:reliability}
\end{table}

\subsection{Provenance: scores a third party can re-derive}\label{sec:provenance}

Every bench emits a \emph{run manifest}: a content hash of the pack (order-
and formatting-independent), the policy label, seed, trial count, library
version, and the per-scenario outcomes, folded into one \texttt{run\_hash}.
\texttt{pack verify-run} re-executes the manifest and reports three booleans:
\texttt{pack\_intact} (the pack on disk still matches the fingerprint),
\texttt{reproduced} (the re-run derived the same \texttt{run\_hash} ---
expected for deterministic policies), and \texttt{within\_tolerance} (for
stochastic models, the pass rate agrees within a stated bound). In the
calibration run, the honest manifest reproduces exactly
($\Delta$ pass rate $0.0$) and a one-word edit to a task after the fact flips
\texttt{pack\_intact} to false --- the pack cannot be quietly reshaped under a
published score. The public leaderboard stores \texttt{run\_hash} and pack
fingerprint with each submission, rejects manifests whose pass rate
contradicts their own rows, and marks reproducible entries; \texttt{agentsynth
diff} compares any two manifests and names the scenarios that regressed,
which is the shape a CI gate wants.

\section{Evaluating agents as-is}\label{sec:interfaces}

Verification is only useful if real agents can reach it, and most agent
evaluation still assumes the agent will be rewritten inside the harness's
framework. \agentsynth{} inverts this: the harness speaks to the agent.

\textbf{One payload, three transports.} Each step, the harness sends one JSON
object --- task, tool catalog (name / description / JSON-schema parameters),
the newest observation, the step budget, and a transcript --- and reads one
action back: \texttt{\{"tool": ..., "args": ...\}} or
\texttt{\{"answer": ...\}}.

\begin{lstlisting}
agentsynth bench --pack core_v2 --agent "python my_agent.py"  # stdio
agentsynth bench --pack core_v2 --agent http://host:8088/act  # HTTP
agentsynth serve-mcp --pack core_v1                           # MCP
\end{lstlisting}

The stdio transport is one JSON line in, one out; a script that exits after
each reply is respawned per step, replies that fail to parse grade as final
answers rather than crashing the run, and a timeout bounds each step. The MCP
mode inverts control entirely: the pack becomes an MCP server
(\texttt{current\_task}, the environment's real tools,
\texttt{submit\_answer}), so an unmodified MCP client --- Claude Code, Claude
Desktop, any agent speaking the protocol~\cite{mcp} --- can be benched by
pointing it at the command. Transport overhead is negligible: on scripted
episodes the stdio protocol adds ${\sim}0.1$\,ms per episode over an
in-process policy, including the process spawn amortized across a five-episode
suite.

\textbf{Exports.} A scenario's checkers \emph{are} a verifiable reward
($1.0$ = goal state reached), so packs export mechanically:
\texttt{pack export --format verifiers} emits an environment module for the
\texttt{verifiers} / Environments Hub ecosystem~\cite{verifiers}, and
\texttt{--format openenv} emits an OpenEnv-style module~\cite{openenv}, both
wrapping the same \texttt{scenario\_reward}. A local-model backend
(\texttt{AGENTSYNTH\_API\_BASE}) drives any OpenAI-compatible server (vLLM,
Ollama) with no provider key, which matters for RL loops that hit the reward
thousands of times.

\section{The synthesis loop}\label{sec:synthesis}

The same machinery generates data. The generator produces multi-turn
trajectories (single-agent tool use, multi-agent hand-offs, grounded code
execution) against the real environments; a six-dimension LLM judge and the
trajectory verifier score them; MinHash dedup, canary insertion, and PII
redaction run at export; and datasets emit as JSONL, ShareGPT, or preference
pairs for SFT / DPO / RFT. Two loops close the circle. \emph{Demo to pack}:
\texttt{pack new --from-demo} turns worked demonstrations into scenarios whose
checkers are derived from the world-diff the demonstration produced --- packs
that validate and audit clean out of the box. \emph{Failure flywheel}: judge a
dataset, mine the failures, and regenerate targeted patches. Provenance
manifests accompany generated datasets as well, so a published dataset pins
the exact pack and policies that produced it.

\section{Limitations}\label{sec:limits}

\textbf{The verifier problem does not disappear.} Outcome checkers exist for
worlds with inspectable state --- databases, files, code under test, API
surfaces. Tasks whose success lives in prose quality still need judges, and
our audit's job there is only to \emph{label} text-graded scenarios honestly
(Table~\ref{tab:audit} shows exactly this weakness in our own refusal tasks).
\textbf{Adversaries are a floor, not a proof.} The audit's knowledge-free
policies catch vacuous and leaky checkers; they do not certify robustness
against a capable adversarial model in the loop, which is the natural next
escalation~\cite{hackerfixer,gamingverifiers}. \textbf{Siblings relabel; they
do not re-plot.} Held-out siblings defeat surface memorization, not
memorization of task \emph{structure}; deeper isomorphisms (schema renames,
value re-plots that preserve orderings) are future work.
\textbf{Scale.} The shipped packs are seeds (4--14 scenarios each),
deliberately small enough to audit by hand while the trust tooling matured;
scaling the pack library without diluting the audit gate is the current
workstream. \textbf{Judged dimensions are excluded from benchmark scores}, by
design; nothing here validates LLM-as-judge quality.

\section{Conclusion}

The Berkeley RDI results made explicit what single-run, transcript-graded
leaderboards had been hiding: most agent benchmarks measure what an agent
\emph{says}, and what an agent says can be gamed. The alternative is not a
bigger benchmark but a stricter contract --- score the world, attack your own
checkers before shipping them, publish scores a stranger can re-derive, and
make honest reporting (decay curves, audit floors, contamination probes) the
default output of the tooling rather than an act of discipline. \agentsynth{}
is one small, complete implementation of that contract, open source and
offline-first, with every claim in this report regenerated by a single script.
We would like the ``Verified'' label to mean something; this is our proposal
for what it should mean.

\paragraph{Availability.} PyPI \texttt{agentsynth-ai} (MIT);
source and packs at \url{https://github.com/agentsynth/agentsynth};
live playground and leaderboard at \url{https://agentsynth.tech}.
\texttt{python paper/experiments.py} regenerates every number in this report
offline in a few minutes on a laptop.

\bibliographystyle{plain}
\bibliography{references}

\appendix

\section{The step payload}\label{app:payload}

The complete integration surface for an external agent, as one stdio exchange:

\begin{lstlisting}
-> {"task": "Refund order 7 ...", "observation": "...",
    "step": 2, "max_steps": 8,
    "tools": [{"name": "sql_query", "description": "...",
               "parameters": {"type": "object", ...}}],
    "transcript": "call sql_query({...})\n-> status\npaid\n(1 row)"}
<- {"tool": "sql_query",
    "args": {"query": "UPDATE orders SET status='refunded' WHERE id=7"}}
...
<- {"answer": "refund issued for order 7"}
\end{lstlisting}

\section{Reproducing this report}\label{app:repro}

\begin{lstlisting}
pip install "agentsynth-ai[dev]"
git clone https://github.com/agentsynth/agentsynth && cd agentsynth
python paper/experiments.py     # writes paper/numbers.json
\end{lstlisting}

The script runs the audits (Table~\ref{tab:audit}), the policy grid
(Table~\ref{tab:bench}), the sibling experiment (Table~\ref{tab:siblings}),
the reliability decay (Table~\ref{tab:reliability}), the provenance
calibration (\S\ref{sec:provenance}), and the transport overhead measurement
(\S\ref{sec:interfaces}), entirely offline with scripted policies and seeded
worlds.

\end{document}
